% Blueprint: the order-zero Riemann--Siegel formula for Hardy's Z.
%
% References: \cite{Siegel1932}; \cite{Edwards} Ch.~7; \cite{Gabcke},
% \cite{AriasDeReyna2011} for the sharp constants at higher order.

\chapter{The order-zero Riemann--Siegel formula}
\label{cha:rs_zero}

Throughout this chapter $t>0$ is real,
\[
  s:=\tfrac12+it,\qquad T:=\frac{t}{2\pi},\qquad \nu:=\lfloor\sqrt T\rfloor ,
\]
$\theta$ is the Riemann--Siegel theta function of \Cref{def:nf_theta}, $Z$ is Hardy's
function of \Cref{def:hz_Z}, and $I(s,a)$ is the contour integral of
\Cref{def:rs_integral}.  We write $\overline{z}$ for complex conjugation.

The proved chain \Cref{thm:rsz_reduction} $\to$ \Cref{thm:rsz_formula} $\to$
\Cref{cor:rsz_bound} uses exactly two facts from \Cref{cha:riemann_siegel}:
\Cref{cor:rs_partial_sum}, which says
\[
  I\bigl(s,\nu+\tfrac12\bigr)-I\bigl(s,\tfrac12\bigr)=\sum_{n=1}^{\nu}n^{-s},
\]
and \Cref{cor:rs_remainder_explicit}, which bounds $I(s,\nu+\tfrac12)$ by $95\,T^{-1/4}$.
\Cref{sec:rsz_siegel}, which proves Siegel's identity, draws on considerably more of
that chapter: the kernel and its residues (\Cref{def:rs_kernel},
\Cref{lem:rs_kernel_residue}), the contour and the branch (\Cref{def:rs_line},
\Cref{def:rs_power}, \Cref{lem:rs_branch_ok}), the constancy of
$\operatorname{Re}(e^{i\pi/4}x)$ along $L(a)$ (\Cref{lem:rs_rotate}), the Gaussian modulus
(\Cref{lem:rs_gauss_modulus}), integrability (\Cref{prop:rs_integrable}), the whole rectangle
argument of \Cref{thm:rs_shift}, and the elementary identity of \Cref{rem:rs_numerics}.

\section{The main sum and the exact reduction}
\label{sec:rsz_reduction}

\begin{definition}[The Riemann--Siegel main sum]
\label{def:rsz_main_sum}
\lean{Zeta.RiemannSiegel.mainSum}
\leanok
\uses{def:nf_theta}
For $t\in\mathbb{R}$ and $\nu\in\mathbb{N}$ set
\[
  \mathcal{S}(t,\nu)\;:=\;2\sum_{n=1}^{\nu}\frac{\cos\bigl(\theta(t)-t\log n\bigr)}{\sqrt n}
  \;\in\;\mathbb{R}.
\]
\end{definition}

The following identity is the entire content of the passage from the contour to the main
sum.  It is unconditional, it holds for every $\nu$, and it is exact: no error term has
been discarded.

\begin{theorem}[Exact reduction]
\label{thm:rsz_reduction}
\lean{Zeta.RiemannSiegel.mainSum_eq_two_re_exp_mul_contourIntegral_sub}
\leanok
\uses{def:rsz_main_sum, cor:rs_partial_sum, def:rs_integral}
For every $t\in\mathbb{R}$ and every $\nu\in\mathbb{N}$, with $s=\tfrac12+it$,
\[
  \mathcal{S}(t,\nu)\;=\;
  2\operatorname{Re}\Bigl(e^{i\theta(t)}\,I\bigl(s,\nu+\tfrac12\bigr)\Bigr)
  -2\operatorname{Re}\Bigl(e^{i\theta(t)}\,I\bigl(s,\tfrac12\bigr)\Bigr).
\]
\end{theorem}

\begin{proof}
\leanok
\uses{cor:rs_partial_sum}
By \Cref{cor:rs_partial_sum},
$I(s,\nu+\tfrac12)-I(s,\tfrac12)=\sum_{n=1}^{\nu}n^{-s}$.  Multiplying by
$e^{i\theta(t)}$ and taking twice the real part, it suffices to check that for every
integer $n\ge1$
\[
  2\operatorname{Re}\bigl(e^{i\theta(t)}n^{-s}\bigr)
  =\frac{2\cos\bigl(\theta(t)-t\log n\bigr)}{\sqrt n}.
\]
Since $n$ is a positive real, the principal power satisfies
$n^{-s}=e^{-s\log n}=n^{-1/2}e^{-it\log n}$, so
$e^{i\theta(t)}n^{-s}=n^{-1/2}e^{i(\theta(t)-t\log n)}$, whose real part is
$n^{-1/2}\cos(\theta(t)-t\log n)$.
\end{proof}

\begin{remark}[The residual term]
\label{rem:rsz_left}
\Cref{thm:rsz_reduction} isolates the whole difficulty.  For $t\ge32\pi$ and the
particular choice $\nu=\lfloor\sqrt T\rfloor$, the term
$2\operatorname{Re}(e^{i\theta}I(s,\nu+\tfrac12))$ is bounded by $2\cdot95\,T^{-1/4}$ by
\Cref{cor:rs_remainder_explicit}; the term $2\operatorname{Re}(e^{i\theta}I(s,\tfrac12))$
is independent of $\nu$ and is \emph{not} small: it is, up to sign, all of $Z(t)$.
Identifying it is \Cref{thm:rsz_siegel}.
\end{remark}

\section{Siegel's identity}
\label{sec:rsz_siegel}

The residual $I(s,\tfrac12)$ left by \Cref{thm:rsz_reduction} is, up to sign, all of
$\zeta(s)$.  Identifying it is Siegel's theorem \cite{Siegel1932} (see also \cite{Edwards},
\S7.9):
\[
  \zeta\bigl(\tfrac12+it\bigr)\;=\;-\,I\bigl(s,\tfrac12\bigr)
     \;-\;e^{-2i\theta(t)}\,\overline{I\bigl(s,\tfrac12\bigr)} ,
\]
which is \Cref{thm:rsz_siegel}, stated and proved at the end of this section.  The section
first records the equivalent general-$s$ form and the shape of the difficulty
(\Cref{rem:rsz_general}--\Cref{rem:rsz_cancellation}), and then carries out Riemann's proof
in three steps: the closed form of an auxiliary Mordell integral
(\Cref{sec:rsz_mordell}), a Mellin transform that crosses the strip
(\Cref{sec:rsz_mellin}), and the assembly (\Cref{sec:rsz_proof}).

\begin{remark}[Equivalent general-$s$ form: the approximate functional equation]
\label{rem:rsz_general}
\uses{thm:rsz_siegel, cor:rs_partial_sum}
Let $\chi(s)=2^{s}\pi^{s-1}\sin\bigl(\tfrac{\pi s}2\bigr)\Gamma(1-s)$ be the factor of the
functional equation, so that $\zeta(s)=\chi(s)\zeta(1-s)$, and let
\[
  \widetilde I(s,a)\;:=\;\overline{I\bigl(\overline{s},a\bigr)}
\]
be the \emph{mirror} contour integral: the integral of $x^{-s}\widetilde G(x)$, with
$\widetilde G(x)=e^{-i\pi x^{2}}/(e^{-i\pi x}-e^{i\pi x})$, over the line of slope $-1$
through $a$, this being exactly the complex conjugate of the $I$-integral by
$x\mapsto\overline x$.  The general form of Siegel's identity is
\begin{equation}
  \zeta(s)+I\bigl(s,\tfrac12\bigr)+\chi(s)\,\widetilde I\bigl(1-s,\tfrac12\bigr)=0
  \qquad(s\in\mathbb{C}\setminus\{1\}).
  \label{eq:rsz_siegel_general}
\end{equation}
Conjugating \Cref{cor:rs_partial_sum} at $1-\overline s$ gives the mirror shift identity
$\sum_{n=1}^{\nu}n^{s-1}=\widetilde I(1-s,\nu+\tfrac12)-\widetilde I(1-s,\tfrac12)$, and
substituting both shift identities into \eqref{eq:rsz_siegel_general} yields the
approximate functional equation in its exact form,
\[
  \zeta(s)=\sum_{n\le\nu}n^{-s}+\chi(s)\sum_{n\le\nu}n^{s-1}
     -I\bigl(s,\nu+\tfrac12\bigr)-\chi(s)\,\widetilde I\bigl(1-s,\nu+\tfrac12\bigr).
\]
On the critical line $1-\overline s=s$, so $\widetilde I(1-s,\tfrac12)=\overline{I(s,\tfrac12)}$,
and $\chi(\tfrac12+it)=e^{-2i\theta(t)}$; thus \eqref{eq:rsz_siegel_general} specialises
to \Cref{thm:rsz_siegel}.  Conversely \eqref{eq:rsz_siegel_general} follows from its
restriction to the critical line only by analytic continuation, which is why the
higher-order chapters use the general form.
Two cautions about the shape of \eqref{eq:rsz_siegel_general}.  First, $\chi$ has poles:
the poles of $\Gamma(1-s)$ at $s=1,2,3,\dots$ are cancelled by the zeros of
$\sin(\pi s/2)$ only at the \emph{even} integers, so $\chi$ is meromorphic with simple
poles at $s=1,3,5,\dots$.  At those odd $s\ge3$ the product
$\chi(s)\widetilde I(1-s,\tfrac12)$ is of the form $\infty\cdot0$ (the identity forces
$\widetilde I(1-s,\tfrac12)$ to vanish there), so \eqref{eq:rsz_siegel_general} is
read on $\mathbb{C}\setminus\{1,3,5,\dots\}$ and extended to the odd integers by
continuity.  \Cref{rem:rsz_hankel_mirror} rewrites the offending product in a form,
$(2\pi)^{s}/2\Gamma(s)$ times an integral, that is free of these spurious poles.  Secondly,
the left-hand side of
\eqref{eq:rsz_siegel_general} is meromorphic on all of $\mathbb{C}$ (the
map $s\mapsto I(s,a)$ is entire for each fixed $a$, and $\zeta$ and $\chi$ are
meromorphic) and the right-hand side is $0$; since two meromorphic functions on the
connected domain $\mathbb{C}$ that agree on a nonempty open set agree everywhere, it
suffices to prove \eqref{eq:rsz_siegel_general} on any nonempty open set.
\end{remark}

\begin{remark}[Sign check]
\label{rem:rsz_signcheck}
The sign in \eqref{eq:rsz_siegel_general} is forced: it is the unique one for which the
displayed approximate functional equation has the classical shape
$\zeta=\sum n^{-s}+\chi(s)\sum n^{s-1}+\text{small}$.  Equivalently, writing
$K(s):=\lim_{a\to\infty}I(s,a)$ along $a\in\tfrac12+\mathbb{Z}_{\ge0}$ (a limit which
exists for $\operatorname{Re}s>1$ because $I(s,\nu+\tfrac12)=\sum_{n\le\nu}n^{-s}+I(s,\tfrac12)$
by \Cref{cor:rs_partial_sum}), the identity \eqref{eq:rsz_siegel_general} reads
$K(s)=-\chi(s)\widetilde I(1-s,\tfrac12)$, and, after one application of the functional
equation, also $I(s,\tfrac12)=-\chi(s)\widetilde K(1-s)$.
\end{remark}

\begin{lemma}[The far-contour limit]
\label{lem:rsz_limit}
\lean{Zeta.RiemannSiegel.tendsto_contourIntegral_atTop_nhds_riemannZeta_add_contourIntegral_half}
\leanok
\uses{def:rs_integral, cor:rs_partial_sum}
For $\operatorname{Re}s>1$ the limit $K(s)$ of \Cref{rem:rsz_signcheck} exists, and
\[
  \lim_{\nu\to\infty}I\bigl(s,\nu+\tfrac12\bigr)\;=\;\zeta(s)+I\bigl(s,\tfrac12\bigr).
\]
\end{lemma}

\begin{proof}
\leanok
\uses{cor:rs_partial_sum}
By \Cref{cor:rs_partial_sum}, $I(s,\nu+\tfrac12)=\sum_{n=1}^{\nu}n^{-s}+I(s,\tfrac12)$ for every
$\nu\ge0$.  For $\operatorname{Re}s>1$ the Dirichlet series converges to $\zeta(s)$, so its
partial sums tend to $\zeta(s)$; adding the constant $I(s,\tfrac12)$ preserves the limit.
\end{proof}

Together with \Cref{rem:rsz_signcheck} this puts \Cref{thm:rsz_siegel}, on the half-plane
$\operatorname{Re}s>1$, in the equivalent form $K(s)=-\chi(s)\widetilde I(1-s,\tfrac12)$: the
content of the identity is the \emph{evaluation of the far-contour limit}, stated in terms of
$I$ alone, with no deformation and no Hankel ray.  This is what \Cref{prop:rsz_general_half}
establishes; its proof evaluates $\zeta(s)+I(s,\tfrac12)$ directly, by \Cref{thm:rsz_mellin}.
\Cref{lem:rsz_limit} is what makes \eqref{eq:rsz_siegel_general} recognisable as an
approximate functional equation, and \Cref{rem:rsz_cancellation} refers to it.

\subsection{The Hankel-ray representation}

The first ingredient collapses the two-dimensional contour integral $I(s,\tfrac12)$ into a
single integral over the positive reals, with the factor $\sin(\pi s/2)$ of $\chi(s)$
already visible.  It is the analogue, for the Riemann--Siegel kernel, of the passage from
Riemann's Hankel loop to $\Gamma(1-s)\zeta(s)$.

It is convenient to prove it in two independent halves.  The contour $L(\tfrac12)$ is first
moved onto the diagonal $L(0)$ (\Cref{lem:rsz_deform}), and the two Hankel rays that make up
$L(0)$ are then evaluated directly (\Cref{lem:rsz_rays}).  This replaces the classical
indented contour: the hypothesis $\sigma<0$ makes $x^{-s}G(x)$ absolutely integrable
\emph{through} the pole of $G$ at the origin, so no arc need be excised, and the deformation
can be run as a limit of contours that each avoid the pole.

\begin{definition}[The Hankel ray integral]
\label{def:rsz_phi}
\lean{Zeta.RiemannSiegel.hankelIntegral}
\leanok
For $s\in\mathbb{C}$ set
\[
  \Phi(s)\;:=\;\int_{0}^{\infty}\frac{\rho^{-s}\,e^{-\pi\rho^{2}}}
       {\sin\bigl(\pi\rho\,e^{i\pi/4}\bigr)}\,d\rho .
\]
\end{definition}

\begin{lemma}[Convergence of $\Phi$]
\label{lem:rsz_phi_conv}
\lean{Zeta.RiemannSiegel.integrableOn_hankelIntegrand}
\leanok
\uses{def:rsz_phi}
For $s=\sigma+i\tau$ with $\sigma<0$ the integral $\Phi(s)$ converges absolutely.
\end{lemma}

\begin{proof}
\leanok
For $\rho>0$ the point $\rho e^{i\pi/4}$ has imaginary part $\rho/\sqrt2>0$, hence is not an
integer, so the integrand is continuous on $(0,\infty)$.  By the identity
$\lvert\sin(u+iv)\rvert^{2}=\sin^{2}u+\sinh^{2}v$ of \Cref{rem:rs_numerics},
\[
  \bigl\lvert\sin(\pi\rho e^{i\pi/4})\bigr\rvert\;\ge\;\sinh(\pi\rho/\sqrt2)\;\ge\;\pi\rho/\sqrt2 ,
\]
so the integrand is $O\bigl(\rho^{-\sigma-1}e^{-\pi\rho^{2}}\bigr)$.  That majorant is
integrable on $(0,\infty)$: at the origin because $-\sigma-1>-1$, which is exactly the
hypothesis $\sigma<0$, and at infinity because of the Gaussian.
\end{proof}

\begin{lemma}[The two Hankel rays]
\label{lem:rsz_rays}
\lean{Zeta.RiemannSiegel.contourIntegral_zero_eq_neg_exp_mul_sin_mul_hankelIntegral}
\leanok
\uses{def:rsz_phi, lem:rsz_phi_conv, def:rs_integral, def:rs_kernel, def:rs_line, def:rs_power}
Let $\sigma<0$.  Then $x\mapsto x^{-s}G(x)$ is absolutely integrable along $L(0)$, and
\[
  I(s,0)\;=\;-\,e^{i\pi(1+s)/4}\,\sin\Bigl(\frac{\pi s}{2}\Bigr)\,\Phi(s).
\]
\end{lemma}

\begin{proof}
\leanok
\uses{def:rs_line, def:rs_power, lem:rsz_phi_conv}
By \Cref{def:rs_line}, $L(0)(w)=\sqrt2\,w\,e^{i\pi/4}$, so the contour through the origin is
the union of the outgoing ray $\{\rho e^{i\pi/4}:\rho>0\}$, traversed for $w>0$, and the
incoming ray $\{\rho e^{-3i\pi/4}:\rho>0\}$, traversed for $w<0$; in both cases
$\rho=\sqrt2\lvert w\rvert$.

On both rays $x^{2}=i\rho^{2}$, so the Gaussian factor of the kernel is
$e^{i\pi x^{2}}=e^{-\pi\rho^{2}}$, and the denominator of $G$ is $2i\sin(\pi\rho e^{i\pi/4})$ on
the outgoing ray; on the incoming ray it is its negative, because $G(-x)=-G(x)$,
immediate from \Cref{def:rs_kernel}.  For the power, \Cref{def:rs_power} gives
$\arg_{\Omega}(\rho e^{i\pi/4})=\tfrac\pi4$ and $\arg_{\Omega}(\rho e^{-3i\pi/4})=-\tfrac{3\pi}4$,
whence
\[
  (\rho e^{i\pi/4})^{-s}=\rho^{-s}e^{-i\pi s/4},\qquad
  (\rho e^{-3i\pi/4})^{-s}=\rho^{-s}e^{3i\pi s/4}.
\]
Thus the integrand along $L(0)$ is $\rho^{-s}e^{-\pi\rho^{2}}/\bigl(2i\sin(\pi\rho e^{i\pi/4})\bigr)$
times the constant $e^{-i\pi s/4}$ on the outgoing ray and times $-e^{3i\pi s/4}$ on the incoming
one; in particular its modulus is a constant multiple of that of the integrand of $\Phi$, so
absolute integrability along $L(0)$ is \Cref{lem:rsz_phi_conv} in the substitution
$\rho=\sqrt2\lvert w\rvert$, which contributes the Jacobian $\sqrt2$.

Collecting the two rays and the velocity $1+i=\sqrt2\,e^{i\pi/4}$ of the parametrisation, the two
factors $\sqrt2$ cancel and
\[
  I(s,0)=e^{i\pi/4}\,\frac{e^{-i\pi s/4}-e^{3i\pi s/4}}{2i}\,\Phi(s)
       =-\,e^{i\pi(1+s)/4}\sin\Bigl(\frac{\pi s}{2}\Bigr)\Phi(s),
\]
using $e^{-i\pi s/4}-e^{3i\pi s/4}=e^{i\pi s/4}\bigl(e^{-i\pi s/2}-e^{i\pi s/2}\bigr)
=-2i\,e^{i\pi s/4}\sin(\pi s/2)$.
\end{proof}

The deformation onto $L(0)$ needs one estimate that \Cref{cha:riemann_siegel} does not supply.
Its bound \Cref{lem:rs_sin_lower}, $\lvert\sin\pi x\rvert\ge\tfrac{\sqrt2}2$ along $L(a)$, holds only for half-integer $a$
and is false as $a\to0$, since $L(0)$ passes through the zero of $\sin(\pi\cdot)$ at the origin.

\begin{lemma}[Lower bound for the kernel denominator on the near contours]
\label{lem:rsz_sin_lb}
\lean{Zeta.RiemannSiegel.norm_sin_pi_mul_line_ge_min}
\leanok
\uses{def:rs_line}
For $0\le a\le\tfrac12$ and every $w\in\mathbb{R}$, writing $x=L(a)(w)$,
\[
  \lvert\sin(\pi x)\rvert\;\ge\;\min\bigl(\tfrac12,\lvert x\rvert\bigr).
\]
\end{lemma}

\begin{proof}
\leanok
Write $x=u+iv$ with $u=a+w$ and $v=w$, and recall
$\lvert\sin\pi x\rvert^{2}=\sin^{2}(\pi u)+\sinh^{2}(\pi v)$ (\Cref{rem:rs_numerics}), so that
$\lvert\sin\pi x\rvert\ge\pi\lvert v\rvert$ and $\lvert\sin\pi x\rvert\ge\lvert\sin\pi u\rvert$.

If $\lvert v\rvert\ge\tfrac14$ then $\lvert\sin\pi x\rvert\ge\pi/4>\tfrac12$ and there is nothing
more to prove.  So assume $\lvert v\rvert<\tfrac14$; then $u=a+v\in(-\tfrac14,\tfrac34)$.  If
$u\le\tfrac12$ then $\lvert u\rvert\le\tfrac12$, and Jordan's inequality gives
$\lvert\sin\pi u\rvert\ge2\lvert u\rvert$; combined with
$\lvert\sin\pi x\rvert\ge\pi\lvert v\rvert\ge2\lvert v\rvert$ this yields
$2\lvert\sin\pi x\rvert^{2}\ge4u^{2}+4v^{2}=4\lvert x\rvert^{2}$, so
$\lvert\sin\pi x\rvert\ge\sqrt2\,\lvert x\rvert\ge\lvert x\rvert$.  (Note that the two bounds must
be \emph{averaged}, not added: each controls only one of $u,v$, so they give the factor $\sqrt2$
and not $2$.  The factor $1$ is all the statement needs.)
If instead $u>\tfrac12$ then $t:=\pi u-\tfrac\pi2\in(0,\tfrac\pi4)$ and
$\sin(\pi u)=\cos t\ge1-\tfrac{t^{2}}2\ge1-\tfrac{\pi^{2}}{32}\ge\tfrac12$, using $\pi<4$.
\end{proof}

\begin{lemma}[Moving the contour onto the diagonal]
\label{lem:rsz_deform}
\lean{Zeta.RiemannSiegel.contourIntegral_zero_eq_contourIntegral_half}
\leanok
\uses{def:rs_integral, def:rs_line, def:rs_power, prop:rs_integrable, lem:rs_kernel_residue, lem:rs_branch_ok, lem:rs_gauss_modulus, lem:rsz_sin_lb, thm:rs_shift}
For $\sigma<0$ one has $I\bigl(s,\tfrac12\bigr)=I(s,0)$.
\end{lemma}

\begin{proof}
\leanok
\uses{prop:rs_integrable, lem:rs_kernel_residue, lem:rs_branch_ok, lem:rs_gauss_modulus, lem:rsz_sin_lb, thm:rs_shift}
\emph{Step 1: all crossing points in $(0,\tfrac12]$ give the same integral.}  Let
$0<b\le a\le\tfrac12$.  First, the integrand is absolutely integrable along $L(a)$ for every such
$a$: by \Cref{lem:rsz_sin_lb} and $\lvert L(a)(w)\rvert\ge a/\sqrt2$ the
kernel denominator is bounded below by the positive constant $\min(\tfrac12,a/\sqrt2)$, and with that
in place the proof of \Cref{prop:rs_integrable} applies verbatim, the constant $\tfrac{\sqrt2}2$ of
the half-integer case being replaced by that constant.

Now argue as in \Cref{thm:rs_shift}.  In the rotated coordinate $z=e^{i\pi/4}x$ the two contours
become the vertical lines $\operatorname{Re}z=b/\sqrt2$ and $\operatorname{Re}z=a/\sqrt2$, and the
rotated integer $e^{i\pi/4}n$ has abscissa $n/\sqrt2$.  No integer lies in $[b,a]$, because
$[b,a]\subseteq(0,\tfrac12]$ and the least positive integer is $1$; hence on the closed rectangle
with those vertical sides and horizontal sides at heights $\pm H$ the integrand is holomorphic:
the branch by \Cref{lem:rs_branch_ok}, since $\operatorname{Re}z\ge b/\sqrt2>0$ there, and the kernel
by \Cref{lem:rs_kernel_residue}.  Cauchy's theorem for a rectangle therefore returns $0$ where
\Cref{thm:rs_shift} returned a residue.  The horizontal sides vanish as $H\to\infty$ by the Gaussian
estimate used there (\Cref{lem:rs_gauss_modulus}), and each vertical side converges to the
corresponding contour integral by dominated convergence, using the integrability just established.
Hence $I(s,a)=I(s,b)$.

\emph{Step 2: the crossing point may be sent to $0$.}  It suffices to exhibit a function
$g\in L^{1}(\mathbb{R})$ with
\[
  \bigl\lvert L(a)(w)^{-s}\,G\bigl(L(a)(w)\bigr)\bigr\rvert\;\le\;g(w)
  \qquad\bigl(0<a\le\tfrac12,\ w\ne0\bigr),
\]
for then, since $L(a)(w)\to L(0)(w)$ as $a\to0$ and the integrand is continuous at $L(0)(w)$ for
$w\ne0$ (there $L(0)(w)$ is neither an integer nor on the branch cut, its rotation
$e^{i\pi/4}L(0)(w)=\sqrt2\,w\,i$ being purely imaginary and nonzero), dominated convergence along
any sequence $a_{n}\downarrow0$ gives $I(s,a_{n})\to I(s,0)$, while $I(s,a_{n})=I(s,\tfrac12)$ for
all $n$ by Step 1.

For the majorant, write $x=L(a)(w)$, so that $\lvert w\rvert\le\lvert x\rvert\le\tfrac12+2\lvert w\rvert$.
By \Cref{lem:rs_branch_ok} the argument of the power contributes at most
$e^{3\pi\lvert\tau\rvert/4}$; by \Cref{lem:rsz_sin_lb},
\[
  \frac1{2\lvert\sin\pi x\rvert}\;\le\;\frac1{2\min(\tfrac12,\lvert x\rvert)}\;\le\;1+\frac1{2\lvert x\rvert};
\]
and the Gaussian factor is $e^{-2\pi(a+w)w}\le e^{-2\pi w^{2}+\pi\lvert w\rvert}$ because
$a\le\tfrac12$.  Hence, with $p:=-\sigma>0$,
\[
  \bigl\lvert x^{-s}G(x)\bigr\rvert\;\le\;e^{3\pi\lvert\tau\rvert/4}
     \Bigl(\lvert x\rvert^{p}+\tfrac12\lvert x\rvert^{p-1}\Bigr)e^{-2\pi w^{2}+\pi\lvert w\rvert}.
\]
Now $\lvert x\rvert^{p}\le e^{p(1/2+2\lvert w\rvert)}$ and
$\lvert x\rvert^{p-1}\le\lvert w\rvert^{p-1}+e^{(p+1)(1/2+2\lvert w\rvert)}$, the first by
$\log t\le t-1$ and $\lvert x\rvert\le\tfrac12+2\lvert w\rvert$, the second by monotonicity of
$t\mapsto t^{p-1}$ (increasing if $p\ge1$, in which case the upper bound on $\lvert x\rvert$ is
used and again converted to an exponential, and decreasing if $p<1$, in which case the lower bound
$\lvert x\rvert\ge\lvert w\rvert$ is used).  Absorbing every linear term into the Gaussian, this
gives a majorant of the shape
\[
  g(w)=C\bigl(1+\lvert w\rvert^{p-1}\bigr)e^{-\pi w^{2}+B\lvert w\rvert},
  \qquad C=C(s),\quad B=2p+2+\pi,
\]
which is integrable: the term without the power is the Gaussian majorant of the proof of
\Cref{thm:rs_remainder_bound}, and the term with it
is integrable at the origin precisely because $p-1>-1$, i.e.\ because $\sigma<0$, and at infinity
because of the Gaussian.  Note that the singular factor cannot be avoided: it is the trace of the
pole of $G$ at the origin, and it is integrable only under the hypothesis $\sigma<0$.
\end{proof}

\begin{lemma}[Hankel-ray representation of $I(s,\tfrac12)$]
\label{lem:rsz_hankel}
\lean{Zeta.RiemannSiegel.contourIntegral_half_eq_neg_exp_mul_sin_mul_hankelIntegral}
\leanok
\uses{lem:rsz_rays, lem:rsz_deform, def:rsz_phi}
Let $s=\sigma+i\tau$ with $\sigma<0$.  Then $\Phi(s)$ converges absolutely and
\[
  I\bigl(s,\tfrac12\bigr)\;=\;-\,e^{i\pi(1+s)/4}\,\sin\Bigl(\frac{\pi s}{2}\Bigr)\,\Phi(s).
\]
\end{lemma}

\begin{proof}
\leanok
\uses{lem:rsz_rays, lem:rsz_deform}
Combine \Cref{lem:rsz_deform} with \Cref{lem:rsz_rays}.
\end{proof}

\begin{remark}[Where the mirror half sits]
\label{rem:rsz_hankel_mirror}
\uses{lem:rsz_hankel, rem:rsz_general}
Conjugating \Cref{lem:rsz_hankel} at $\overline{1-s}$ gives, for $\operatorname{Re}s>1$,
\[
  \widetilde I\bigl(1-s,\tfrac12\bigr)
    =-\,e^{-i\pi(2-s)/4}\cos\Bigl(\frac{\pi s}{2}\Bigr)
      \int_{0}^{\infty}\frac{\rho^{s-1}e^{-\pi\rho^{2}}}{\sin(\pi\rho e^{-i\pi/4})}\,d\rho ,
\]
and, since
$\chi(s)\cos\bigl(\tfrac{\pi s}2\bigr)=2^{s-1}\pi^{s-1}\Gamma(1-s)\sin(\pi s)
=\dfrac{(2\pi)^{s}}{2\,\Gamma(s)}$,
the general identity \eqref{eq:rsz_siegel_general} is, for $\operatorname{Re}s>1$,
equivalent to the concrete Mellin-type statement
\[
  \zeta(s)+I\bigl(s,\tfrac12\bigr)
   =\frac{(2\pi)^{s}e^{-i\pi(2-s)/4}}{2\,\Gamma(s)}
     \int_{0}^{\infty}\frac{\rho^{s-1}e^{-\pi\rho^{2}}}{\sin(\pi\rho e^{-i\pi/4})}\,d\rho .
\]
The two representations are valid on the disjoint half-planes $\operatorname{Re}s<0$ and
$\operatorname{Re}s>1$, and matching them across the strip is the analytic heart of
\Cref{thm:rsz_siegel}.  The displayed Mellin-type statement is
\Cref{thm:rsz_mellin} and the mirror evaluation above is \Cref{lem:rsz_mirror}; the crossing
is carried out in \Cref{sec:rsz_mellin} by a Mellin transform in an auxiliary variable, not
by any further movement of the contour in $x$.
\end{remark}

\subsection{Gaussian structure of the kernel}

The remaining two ingredients are needed again, unchanged, by the saddle-point chapters.

\begin{lemma}[Half-plane Gaussian expansions of $G$]
\label{lem:rsz_expansion}
\lean{Zeta.RiemannSiegel.kernel_eq_tsum_of_im_neg, Zeta.RiemannSiegel.kernel_eq_neg_tsum_of_im_pos}
\leanok
\uses{def:rs_kernel}
For $\operatorname{Im}x<0$,
\[
  G(x)\;=\;e^{-i\pi/4}\sum_{k\ge0}\exp\Bigl(i\pi\bigl(x-\tfrac{2k+1}2\bigr)^{2}\Bigr),
\]
and for $\operatorname{Im}x>0$,
\[
  G(x)\;=\;-\,e^{-i\pi/4}\sum_{k\ge0}\exp\Bigl(i\pi\bigl(x+\tfrac{2k+1}2\bigr)^{2}\Bigr),
\]
both series converging absolutely, and locally uniformly on the respective open
half-planes.
\end{lemma}

\begin{proof}
\leanok
Let $\operatorname{Im}x<0$.  Then $\lvert e^{-2\pi ix}\rvert=e^{2\pi\operatorname{Im}x}<1$,
so from $e^{i\pi x}-e^{-i\pi x}=e^{i\pi x}\bigl(1-e^{-2\pi ix}\bigr)$ and the geometric
series,
\[
  G(x)=e^{i\pi x^{2}-i\pi x}\sum_{k\ge0}e^{-2\pi ikx}
      =\sum_{k\ge0}\exp\bigl(i\pi\bigl(x^{2}-(2k+1)x\bigr)\bigr).
\]
Completing the square, $x^{2}-(2k+1)x=\bigl(x-\tfrac{2k+1}2\bigr)^{2}-\tfrac{(2k+1)^{2}}4$,
so the $k$-th term is $e^{-i\pi(2k+1)^{2}/4}\exp\bigl(i\pi(x-\tfrac{2k+1}2)^{2}\bigr)$.
Now $(2k+1)^{2}=4k(k+1)+1\equiv1\pmod 8$ because $k(k+1)$ is even, so
$e^{-i\pi(2k+1)^{2}/4}=e^{-i\pi/4}$ for every $k$, which gives the first formula.  The
series converges absolutely and locally uniformly because
$\lvert e^{-2\pi ikx}\rvert=e^{2\pi k\operatorname{Im}x}$ is a geometric series with ratio
$<1$, locally uniformly in the open lower half-plane.

For $\operatorname{Im}x>0$ the same computation applies to
$e^{i\pi x}-e^{-i\pi x}=-e^{-i\pi x}\bigl(1-e^{2\pi ix}\bigr)$, with
$\lvert e^{2\pi ix}\rvert<1$; it produces $x^{2}+(2k+1)x$ in the exponent and an overall
minus sign.
\end{proof}

\begin{lemma}[The basic Gaussian contour integral]
\label{lem:rsz_gauss_integral}
\lean{Zeta.RiemannSiegel.integral_exp_pi_mul_I_mul_line_sub_sq}
\leanok
\uses{def:rs_line}
For every real $a$ and every \emph{complex} $c$,
\[
  \int_{L(a)}e^{i\pi(x-c)^{2}}\,dx
  \;=\;(1+i)\int_{\mathbb{R}}\exp\Bigl(i\pi\bigl(L(a)(w)-c\bigr)^{2}\Bigr)dw
  \;=\;e^{i\pi/4} .
\]
In particular the value does not depend on $a$ or $c$.  Only the reality of $a$, which fixes
the contour, is used.
\end{lemma}

\begin{proof}
\leanok
Put $c':=a-c$ and substitute $\rho:=\sqrt2\,w$, so that
$x-c=c'+\rho e^{i\pi/4}$ and $dx=(1+i)\,dw=e^{i\pi/4}d\rho$.  Then
\[
  i\pi(x-c)^{2}=i\pi c'^{2}+2\pi i c'\rho\,e^{i\pi/4}-\pi\rho^{2}
  =i\pi c'^{2}-\pi(\rho-\alpha)^{2}+\pi\alpha^{2},\qquad \alpha:=c'e^{3i\pi/4},
\]
and $\pi\alpha^{2}=\pi c'^{2}e^{3i\pi/2}=-i\pi c'^{2}$, so the two constants cancel
exactly and
\[
  \int_{L(a)}e^{i\pi(x-c)^{2}}dx=e^{i\pi/4}\int_{\mathbb{R}}e^{-\pi(\rho-\alpha)^{2}}d\rho
  =e^{i\pi/4},
\]
the last step because $\int_{\mathbb{R}}e^{-\pi(\rho-\alpha)^{2}}d\rho=1$ for every
\emph{complex} $\alpha$.  This is a contour shift, not translation invariance of Lebesgue
measure, and mathlib supplies it ready-made: \texttt{integral\_cexp\_quadratic}, in
\texttt{Mathlib/Analysis/SpecialFunctions/Gaussian/FourierTransform.lean}, states that
$\int_{\mathbb{R}}e^{bx^{2}+cx+d}\,dx=(\pi/-b)^{1/2}e^{\,d-c^{2}/4b}$ whenever
$\operatorname{Re}b<0$, for arbitrary complex $c,d$.  Applied with $b=-\pi$,
$c=2\pi ic'e^{i\pi/4}$ and $d=i\pi c'^{2}$ it returns $e^{i\pi c'^{2}}e^{-i\pi c'^{2}}=1$
directly, so a formalization needs neither the substitution $\rho=\sqrt2\,w$ nor a separate
appeal to $\int_{\mathbb{R}}e^{-\pi u^{2}}du=1$: it can be run in the variable $w$ of
\Cref{def:rs_line} as it stands.
\end{proof}

\begin{remark}[Why the modulus bound of \Cref{cha:riemann_siegel} is not enough here]
\label{rem:rsz_cancellation}
\uses{lem:rsz_expansion, lem:rsz_gauss_integral}
On $L(a)$ the modulus of $G$ is enormous: by \Cref{lem:rsz_expansion} the term $k=0$ alone
has modulus $e^{\pi a(a-1)/2}$ at the point $w=-a/2$, whereas for $\operatorname{Re}u>1$ the
integral $I(u,a)$ stays bounded as $a\to\infty$, since
$I(u,\nu+\tfrac12)=\sum_{n\le\nu}n^{-u}+I(u,\tfrac12)$ by \Cref{cor:rs_partial_sum} and the
partial sums converge.  (On the critical line itself no such boundedness holds: there
$\sum_{n\le\nu}n^{-1/2-it}$ grows like $\nu^{1/2}$, and it is the \emph{ratio} of that
growth to $e^{\pi a(a-1)/2}$ that is the point.)  \Cref{lem:rsz_gauss_integral} is exactly the
mechanism: each Gaussian in the expansion integrates, along the slope-$1$ line, to the
single constant $e^{i\pi/4}$, the factor $e^{\pi c'^{2}/2}$ from the peak being cancelled by
the oscillation.  Any attempt to bound \Cref{thm:rsz_siegel} by moduli is therefore
hopeless; only the exact Gaussian evaluations can be used.  This is also why
\Cref{cha:riemann_siegel} bounds $I(s,\nu+\tfrac12)$ only under the hypothesis
$T\le2a^{2}$, which keeps the contour at or beyond the saddle.
\end{remark}

\subsection{Riemann's Mordell integral}
\label{sec:rsz_mordell}

The whole proof rests on one auxiliary function of a \emph{complex} parameter $u$: the
contour integral of the kernel against the exponential $e^{-2\pi iux}$ rather than against
the power $x^{-s}$.  Unlike $I(s,a)$ it can be evaluated in closed form by elementary means,
because the exponential, unlike the power, interacts with both of the two structural
symmetries of $G$: the quasi-periodicity $G(x+1)=e^{2\pi ix}G(x)$ and the pole at every
integer.  This is the integral $\varPhi(\tau,u)$ studied by Kronecker and Mordell, at the
parameter $\tau$ that Riemann needed; see \cite{Siegel1932} \S1.

Two of the three ingredients are already available: \Cref{lem:rs_kernel_residue} for the
residue and \Cref{thm:rs_shift} for the contour shift across it.  The third is
\Cref{lem:rsz_gauss_integral}, which is stated for a complex centre precisely for this use.

\begin{definition}[The Mordell integral]
\label{def:rsz_mordell}
\lean{Zeta.RiemannSiegel.mordellIntegral}
\leanok
\uses{def:rs_line, def:rs_kernel}
For $u\in\mathbb{C}$ and $a\in\tfrac12+\mathbb{Z}_{\ge0}$ set
\[
  N(u,a)\;:=\;\int_{L(a)}e^{-2\pi iux}G(x)\,dx
    \;=\;(1+i)\int_{\mathbb{R}}e^{-2\pi iuL(a)(w)}\,G\bigl(L(a)(w)\bigr)\,dw ,
\]
and write $N(u):=N(u,\tfrac12)$.
\end{definition}

\begin{lemma}[Convergence]
\label{lem:rsz_mordell_conv}
\lean{Zeta.RiemannSiegel.integrable_mordellIntegrand}
\leanok
\uses{def:rsz_mordell, lem:rs_gauss_modulus, lem:rs_sin_lower, lem:rs_rotate}
For every $u\in\mathbb{C}$ and every $a\in\tfrac12+\mathbb{Z}_{\ge0}$ the integrand of
$N(u,a)$ is absolutely integrable on $\mathbb{R}$.
\end{lemma}

\begin{proof}
\leanok
\uses{lem:rs_gauss_modulus, lem:rs_sin_lower, lem:rs_rotate}
As in \Cref{prop:rs_integrable}, $\lvert G(L(a)(w))\rvert\le\tfrac1{\sqrt2}e^{\pi a^{2}}
e^{-\pi w^{2}}$, using \Cref{lem:rs_gauss_modulus} for the numerator and
\Cref{lem:rs_sin_lower} for the denominator.  By \Cref{lem:rs_rotate},
$\lvert L(a)(w)\rvert\le a+2\lvert w\rvert$, so
$\lvert e^{-2\pi iuL(a)(w)}\rvert\le e^{2\pi\lvert u\rvert(a+2\lvert w\rvert)}$.  A Gaussian
times an exponential of a linear function is integrable.  Note that, in contrast with
\Cref{prop:rs_integrable}, no branch of a power is involved, so the hypothesis $a>0$ is not
needed here; it appears only because \Cref{thm:rs_shift} is stated for half-integers.
\end{proof}

\begin{lemma}[First difference equation: the Gaussian shift]
\label{lem:rsz_mordell_gauss}
\lean{Zeta.RiemannSiegel.mordellIntegral_sub_mordellIntegral_add_one}
\leanok
\uses{def:rsz_mordell, lem:rsz_gauss_integral, lem:rsz_mordell_conv}
For every $u\in\mathbb{C}$ and $a\in\tfrac12+\mathbb{Z}_{\ge0}$,
\[
  N(u,a)-N(u+1,a)\;=\;e^{i\pi/4}\,e^{-i\pi\left(u+\frac12\right)^{2}} .
\]
In particular the difference does not depend on $a$.
\end{lemma}

\begin{proof}
\leanok
\uses{lem:rsz_gauss_integral, lem:rsz_mordell_conv}
Both integrals converge absolutely by \Cref{lem:rsz_mordell_conv}, so they may be
subtracted under the integral sign:
\[
  N(u,a)-N(u+1,a)=\int_{L(a)}\bigl(1-e^{-2\pi ix}\bigr)e^{-2\pi iux}G(x)\,dx .
\]
The pole of $G$ is cancelled exactly: by \Cref{def:rs_kernel},
\[
  \bigl(1-e^{-2\pi ix}\bigr)G(x)
  =\frac{e^{-i\pi x}\bigl(e^{i\pi x}-e^{-i\pi x}\bigr)e^{i\pi x^{2}}}
        {e^{i\pi x}-e^{-i\pi x}}
  =e^{i\pi x^{2}-i\pi x},
\]
an entire function.  Hence the integrand is $e^{i\pi(x^{2}-(2u+1)x)}$, and completing the
square,
\[
  x^{2}-(2u+1)x=\Bigl(x-\bigl(u+\tfrac12\bigr)\Bigr)^{2}-\bigl(u+\tfrac12\bigr)^{2}.
\]
Now apply \Cref{lem:rsz_gauss_integral} with the complex centre $c=u+\tfrac12$.  (This is
where the complex centre is indispensable: $u$ is the variable of the difference equation
and is complex throughout.)
\end{proof}

\begin{lemma}[Second difference equation: the pole shift]
\label{lem:rsz_mordell_pole}
\lean{Zeta.RiemannSiegel.mordellIntegral_sub_exp_mul_mordellIntegral_add_one}
\leanok
\uses{def:rsz_mordell, thm:rs_shift, lem:rs_kernel_residue, lem:rsz_mordell_conv}
For every $u\in\mathbb{C}$,
\[
  N(u)-e^{2\pi iu}\,N(u+1)\;=\;1 .
\]
\end{lemma}

\begin{proof}
\leanok
\uses{thm:rs_shift, lem:rs_kernel_residue, lem:rsz_mordell_conv}
The comparison is between the contours $L(\tfrac12)$ and $L(\tfrac32)$, which enclose the
single pole $x=1$ of $G$.  The rectangle argument of \Cref{thm:rs_shift} applies with the
integrand $x^{-s}G(x)$ replaced by $f_{u}(x):=e^{-2\pi iux}G(x)$; that argument takes the
integrand as a parameter, and its three hypotheses hold here.
\begin{enumerate}
  \item \emph{Holomorphy off the pole.}  $f_{u}$ is holomorphic on the closed strip
        $\tfrac12\le\operatorname{Re}(e^{i\pi/4}x)\sqrt2\le\tfrac32$ minus $\{1\}$, because
        $e^{-2\pi iux}$ is entire and $G$ is holomorphic off $\mathbb{Z}$
        (\Cref{lem:rs_kernel_residue}).  This is strictly easier than for $x^{-s}G(x)$, where
        the branch of the power had to be controlled as well; in particular no analogue of
        \Cref{lem:rs_branch_ok} is needed.
  \item \emph{Simple pole with a computable residue.}  By \Cref{lem:rs_kernel_residue},
        $\lim_{x\to1}(x-1)G(x)=\tfrac1{2\pi i}$, so
        $\lim_{x\to1}(x-1)f_{u}(x)=e^{-2\pi iu}/2\pi i$, the value of $e^{-2\pi iux}$ at
        $x=1$ times that residue.
  \item \emph{Decay of the horizontal sides.}  By the proof of \Cref{lem:rsz_mordell_conv}
        the integrand is bounded by a Gaussian in $w$ times an exponential of a linear
        function, uniformly for $\operatorname{Re}x$ in the compact interval
        $[\tfrac12,\tfrac32]$, so the horizontal sides tend to $0$.
\end{enumerate}
Hence
\[
  N(u,\tfrac32)-N(u,\tfrac12)\;=\;2\pi i\cdot\frac{e^{-2\pi iu}}{2\pi i}\;=\;e^{-2\pi iu}.
\]
On the other hand $L(\tfrac32)(w)=L(\tfrac12)(w)+1$ for every $w$ (the two contours differ
by the translation $x\mapsto x+1$ in the \emph{same} parameter $w$), and $G(x+1)=e^{2\pi ix}G(x)$, directly from \Cref{def:rs_kernel} since
$e^{i\pi(x+1)^{2}}=-e^{i\pi x^{2}}e^{2\pi ix}$ and $e^{\pm i\pi(x+1)}=-e^{\pm i\pi x}$.
Therefore
\[
  N(u,\tfrac32)=\int_{L(\frac12)}e^{-2\pi iu(x+1)}G(x+1)\,dx
    =e^{-2\pi iu}\int_{L(\frac12)}e^{-2\pi i(u-1)x}G(x)\,dx=e^{-2\pi iu}N(u-1).
\]
Combining, $e^{-2\pi iu}N(u-1)-N(u)=e^{-2\pi iu}$, i.e.\ $N(u-1)=e^{2\pi iu}N(u)+1$;
replacing $u$ by $u+1$ and using $e^{2\pi i(u+1)}=e^{2\pi iu}$ gives the claim.
\end{proof}

\begin{theorem}[Closed form of the Mordell integral]
\label{thm:rsz_mordell}
\lean{Zeta.RiemannSiegel.mordellIntegral_eq}
\leanok
\uses{lem:rsz_mordell_gauss, lem:rsz_mordell_pole}
For every $u\in\mathbb{C}\setminus\mathbb{Z}$,
\[
  N(u)\;=\;\int_{L(\frac12)}e^{-2\pi iux}G(x)\,dx
      \;=\;\frac{e^{-i\pi u^{2}}-e^{-i\pi u}}{e^{i\pi u}-e^{-i\pi u}} .
\]
\end{theorem}

\begin{proof}
\leanok
\uses{lem:rsz_mordell_gauss, lem:rsz_mordell_pole}
Eliminate $N(u+1)$ between the two difference equations.  By
\Cref{lem:rsz_mordell_gauss}, $N(u+1)=N(u)-e^{i\pi/4}e^{-i\pi(u+\frac12)^{2}}$; substituting
into \Cref{lem:rsz_mordell_pole},
\[
  N(u)\bigl(1-e^{2\pi iu}\bigr)\;=\;1-e^{2\pi iu}e^{i\pi/4}e^{-i\pi(u+\frac12)^{2}} .
\]
Since $-i\pi(u+\tfrac12)^{2}=-i\pi u^{2}-i\pi u-\tfrac{i\pi}4$, the second term on the right
is $e^{-i\pi u^{2}}e^{2\pi iu-i\pi u}=e^{-i\pi u^{2}+i\pi u}$, so
\[
  N(u)\bigl(1-e^{2\pi iu}\bigr)\;=\;1-e^{-i\pi u^{2}+i\pi u} .
\]
For $u\notin\mathbb{Z}$ the factor $1-e^{2\pi iu}$ is nonzero; dividing and multiplying
numerator and denominator by $e^{-i\pi u}$ gives the stated form.  (Both sides are in fact
entire, as the numerator vanishes at every integer; the ray of \Cref{sec:rsz_mellin} meets
$\mathbb{Z}$ only at $0$.)
\end{proof}

\begin{remark}[Normalisation]
\label{rem:rsz_mordell_numeric}
\uses{thm:rsz_mordell}
The closed form holds for every $a\in(0,1)$, since the argument uses only that the pole $x=0$
lies to the left of $L(a)$ and $x=1$ to its right.  Literature normalisations of the
Mordell integral differ.  The Gaussian here has parameter $\pi$, whereas the model
integral $M(p)$ of \Cref{cha:rs_first} has parameter $2\pi$: they are different members of
Mordell's family and neither closed form implies the other.
\end{remark}

\subsection{The Mellin transform along the diagonal ray}
\label{sec:rsz_mellin}

The passage from \Cref{thm:rsz_mordell} to $\zeta$ is Riemann's, and it is the step that
crosses the strip.  One multiplies the identity of \Cref{thm:rsz_mordell} by $u^{-s}$ and
integrates $u$ over the ray $e^{-i\pi/4}(0,\infty)$.  Three things then happen at once:
\begin{itemize}
  \item the geometric part of $N(u)$ becomes $\sum_{n\ge1}n^{s-1}=\zeta(1-s)$;
  \item the exponential $e^{-2\pi iux}$ inside the contour integral is a Mellin kernel in
        $u$, and integrating it against $u^{-s}$ produces $\Gamma(1-s)$ times $x^{s-1}$,
        so the contour integral $\int_{L(1/2)}e^{-2\pi iux}G(x)dx$ turns into
        $\int_{L(1/2)}x^{s-1}G(x)\,dx=I(1-s,\tfrac12)$;
  \item the Gaussian part of $N(u)$ becomes a Hankel-ray integral of the mirror kernel.
\end{itemize}
No contour in $x$ is moved.  The ray is forced: it must lie in the lower half-plane for the
geometric series to converge, and $e^{-i\pi/4}$ is the unique direction for which the
resulting exponential rate is positive \emph{uniformly along $L(a)$}, by
\Cref{lem:rs_rotate}, which says that $\operatorname{Re}(e^{i\pi/4}x)=a/\sqrt2$ is
\emph{constant} on $L(a)$.  It is that constancy which makes the double integral absolutely
convergent, hence Fubini applicable.

Throughout this subsection $\omega:=e^{-i\pi/4}$, $s=\sigma+i\tau$ and $\sigma<0$, and for
$F:\mathbb{C}\to\mathbb{C}$ we abbreviate
\[
  \mathcal{R}F(s)\;:=\;\omega\int_{0}^{\infty}(\rho\omega)^{-s}F(\rho\omega)\,d\rho ,
\]
the integral of $u^{-s}F(u)$ over the ray, with $(\rho\omega)^{-s}$ the principal branch
(legitimate since $\arg(\rho\omega)=-\tfrac\pi4$).

\begin{lemma}[The $\Gamma$-integral with a complex rate]
\label{lem:rsz_gamma_cplx}
\lean{Complex.integral_cpow_neg_mul_exp_neg_mul_Ioi}
\leanok
For $b\in\mathbb{C}$ with $\operatorname{Re}b>0$ and $s\in\mathbb{C}$ with
$\operatorname{Re}s<1$,
\[
  \int_{0}^{\infty}\rho^{-s}e^{-b\rho}\,d\rho\;=\;\Gamma(1-s)\,b^{\,s-1},
\]
the power being the principal branch.
\end{lemma}

\begin{proof}
\leanok
Mathlib's \texttt{Complex.integral\_cpow\_mul\_exp\_neg\_mul\_Ioi},
$\int_{(0,\infty)}t^{a-1}e^{-rt}dt=(1/r)^{a}\Gamma(a)$ under $0<\operatorname{Re}a$ and $r>0$,
is exactly this statement at $a=1-s$ (whose hypothesis $0<\operatorname{Re}(1-s)$ is the
$\operatorname{Re}s<1$ above), but only for a \emph{real} rate $r$.  The extension is by the
identity theorem in $b$, not by rotating the ray.  Fix $s$ with $\sigma<1$ and let
$H:=\{b:\operatorname{Re}b>0\}$, an open connected set.  The right-hand side is holomorphic
on $H$.  For the left-hand side, on a compact $K\subset H$ one has
$\operatorname{Re}b\ge\delta>0$ and $\lvert b\rvert\le R$, so
$\lvert\rho^{-s}e^{-b\rho}\rvert\le\rho^{-\sigma}e^{\tau\arg\rho}e^{-\delta\rho}
=\rho^{-\sigma}e^{-\delta\rho}$ and
$\lvert\partial_{b}(\rho^{-s}e^{-b\rho})\rvert=\rho^{1-\sigma}e^{-\rho\operatorname{Re}b}
\le\rho^{1-\sigma}e^{-\delta\rho}$, both
integrable on $(0,\infty)$ because $-\sigma>-1$; differentiation under the integral sign
(mathlib's \texttt{hasDerivAt\_integral\_of\_dominated\_loc\_of\_deriv\_le}) makes the
left-hand side holomorphic on $H$ too.  The two agree on $(0,\infty)\subset H$, a set with a
limit point in $H$, hence on all of $H$.

In Lean the identity theorem is applied at $b=1$, the accumulating points being the reals
$1+\varepsilon/2$; the ball on which the dominating functions are used is
$\lvert b-b_{0}\rvert<\delta:=\tfrac12\operatorname{Re}b_{0}$, on which
$\operatorname{Re}b\ge\delta$.  The half-plane is convex, hence preconnected.  Both the
statement above and the strict generalisation of
\texttt{Complex.integral\_cpow\_mul\_exp\_neg\_mul\_Ioi} that it specialises,
\texttt{Complex.integral\_cpow\_mul\_exp\_neg\_mul\_Ioi'}, live in the \texttt{Complex}
namespace in \texttt{Zeta/Analysis/SpecialFunctions/GammaIntegral.lean}, together with the
absolute convergence \texttt{Complex.integrableOn\_cpow\_mul\_exp\_neg\_mul\_Ioi}.
\end{proof}

\begin{definition}[The mirror Hankel ray integral]
\label{def:rsz_phi_mirror}
\lean{Zeta.RiemannSiegel.mirrorHankelIntegral}
\leanok
\uses{def:rsz_phi}
For $s\in\mathbb{C}$ set
\[
  \widetilde\Phi(s)\;:=\;\int_{0}^{\infty}
     \frac{\rho^{-s}e^{-\pi\rho^{2}}}{\sin\bigl(\pi\rho e^{-i\pi/4}\bigr)}\,d\rho
   \;=\;\overline{\Phi(\overline{s})} ,
\]
the complex conjugate of \Cref{def:rsz_phi}.  It converges absolutely for
$\operatorname{Re}s<0$, by \Cref{lem:rsz_phi_conv} and conjugation.  It is the second equality
that is taken as the Lean definition, exactly as for $\widetilde I$ in \Cref{def:rsz_mirror}:
defining the mirror by conjugation transports everything proved about $\Phi$ to
$\widetilde\Phi$.
\end{definition}

\begin{lemma}[The geometric term]
\label{lem:rsz_ray_zeta}
\lean{Zeta.RiemannSiegel.rayIntegral_inv_one_sub_exp}
\leanok
\uses{lem:rsz_gamma_cplx}
For $\sigma<0$,
\[
  \mathcal{R}\Bigl(u\mapsto\frac1{1-e^{2\pi iu}}\Bigr)(s)
    \;=\;-\,(2\pi)^{s-1}e^{i\pi(s-1)/2}\,\Gamma(1-s)\,\zeta(1-s).
\]
\end{lemma}

\begin{proof}
\leanok
\uses{lem:rsz_gamma_cplx}
On the ray, $\operatorname{Im}(\rho\omega)=-\rho/\sqrt2<0$, so
$\lvert e^{2\pi i\rho\omega}\rvert=e^{\sqrt2\pi\rho}>1$ and
\[
  \frac1{1-e^{2\pi iu}}=\frac{-e^{-2\pi iu}}{1-e^{-2\pi iu}}
    =-\sum_{n\ge1}e^{-2\pi inu},\qquad u=\rho\omega,
\]
with $\lvert e^{-2\pi in\rho\omega}\rvert=e^{-2\pi n\rho\operatorname{Re}(i\omega)}
=e^{-\sqrt2\pi n\rho}$, since $i\omega=e^{i\pi/4}$.  The series may be integrated term by
term because
$\sum_{n\ge1}\int_{0}^{\infty}\rho^{-\sigma}e^{-\sqrt2\pi n\rho}d\rho
=\Gamma(1-\sigma)\sum_{n\ge1}(\sqrt2\pi n)^{\sigma-1}<\infty$, the last sum converging
precisely because $1-\sigma>1$.  By \Cref{lem:rsz_gamma_cplx} with
$b=2\pi in\omega=2\pi ne^{i\pi/4}$, whose real part $\sqrt2\pi n$ is positive,
\[
  \omega\int_{0}^{\infty}(\rho\omega)^{-s}e^{-2\pi in\rho\omega}d\rho
   =\omega^{1-s}\Gamma(1-s)\bigl(2\pi ne^{i\pi/4}\bigr)^{s-1}
   =(2\pi n)^{s-1}e^{i\pi(s-1)/2}\Gamma(1-s),
\]
using $(\rho\omega)^{-s}=\rho^{-s}\omega^{-s}$ (valid as $\rho>0$) and
$\omega^{1-s}\bigl(e^{i\pi/4}\bigr)^{s-1}=e^{i\pi(s-1)/4}e^{i\pi(s-1)/4}=e^{i\pi(s-1)/2}$.
Summing over $n$ and using $\sum_{n\ge1}n^{s-1}=\zeta(1-s)$, legitimate because
$\operatorname{Re}(1-s)>1$, gives the claim with the overall minus sign of the geometric
expansion.
\end{proof}

\begin{lemma}[The contour term]
\label{lem:rsz_ray_contour}
\lean{Zeta.RiemannSiegel.rayIntegral_mordellIntegral}
\leanok
\uses{def:rsz_mordell, lem:rsz_gamma_cplx, lem:rs_rotate, def:rs_power, def:rs_integral,
prop:rs_integrable}
For $\sigma<0$,
\[
  \mathcal{R}\bigl(u\mapsto N(u)\bigr)(s)
    \;=\;(2\pi)^{s-1}e^{i\pi(s-1)/2}\,\Gamma(1-s)\,I\bigl(1-s,\tfrac12\bigr).
\]
\end{lemma}

\begin{proof}
\leanok
\uses{lem:rsz_gamma_cplx, lem:rs_rotate, def:rs_power, prop:rs_integrable}
Write out $N$ and apply Fubini on $(0,\infty)\times\mathbb{R}$.  Absolute convergence is the
point, and it is where the choice of ray pays: with $x=L(\tfrac12)(w)$,
\[
  \bigl\lvert e^{-2\pi i\rho\omega x}\bigr\rvert
   =e^{-2\pi\rho\operatorname{Re}(i\omega x)}=e^{-2\pi\rho\operatorname{Re}(e^{i\pi/4}x)}
   =e^{-\pi\rho/\sqrt2},
\]
\emph{independently of $w$}, by \Cref{lem:rs_rotate}.  Together with
$\lvert(\rho\omega)^{-s}\rvert=\rho^{-\sigma}e^{-\pi\tau/4}$ and the Gaussian bound
$\lvert G(x)\rvert\le\tfrac1{\sqrt2}e^{\pi/4}e^{-\pi w^{2}}$ of \Cref{prop:rs_integrable},
the modulus of the double integrand is at most a constant times
$\rho^{-\sigma}e^{-\pi\rho/\sqrt2}e^{-\pi w^{2}}$, which is integrable since $-\sigma>-1$.

For fixed $x$ the inner integral is \Cref{lem:rsz_gamma_cplx} with
$b=2\pi i\omega x=2\pi e^{i\pi/4}x$, whose real part is $2\pi\cdot\tfrac1{2\sqrt2}>0$:
\[
  \omega\int_{0}^{\infty}(\rho\omega)^{-s}e^{-2\pi i\rho\omega x}d\rho
   =\omega^{1-s}\Gamma(1-s)\bigl(2\pi e^{i\pi/4}x\bigr)^{s-1}
   =(2\pi)^{s-1}e^{i\pi(s-1)/4}\Gamma(1-s)\,\omega^{1-s}\,x^{s-1},
\]
where $\bigl(e^{i\pi/4}x\bigr)^{s-1}=e^{i\pi(s-1)/4}x^{s-1}$ is exactly the defining relation
of the branch of \Cref{def:rs_power}.  This is the one place where the slanted branch
of \Cref{def:rs_power} is needed: it is the branch in which the identity holds without a
correction.  Since
$\omega^{1-s}=e^{i\pi(s-1)/4}$ the prefactor is $(2\pi)^{s-1}e^{i\pi(s-1)/2}\Gamma(1-s)$, and
$\int_{L(1/2)}x^{s-1}G(x)dx=I(1-s,\tfrac12)$ by \Cref{def:rs_integral}.
\end{proof}

\begin{lemma}[The Gaussian term]
\label{lem:rsz_ray_hankel}
\lean{Zeta.RiemannSiegel.rayIntegral_mirrorKernel}
\leanok
\uses{def:rsz_phi_mirror}
For $\sigma<0$,
\[
  \mathcal{R}\Bigl(u\mapsto\frac{e^{-i\pi u^{2}}}{e^{i\pi u}-e^{-i\pi u}}\Bigr)(s)
   \;=\;\frac{e^{i\pi(s-1)/4}}{2i}\,\widetilde\Phi(s).
\]
\end{lemma}

\begin{proof}
\leanok
\uses{def:rsz_phi_mirror}
Substitute $u=\rho\omega$.  Then $u^{2}=-i\rho^{2}$, so $e^{-i\pi u^{2}}=e^{-\pi\rho^{2}}$,
and $e^{i\pi u}-e^{-i\pi u}=2i\sin(\pi\rho\omega)$.  With
$(\rho\omega)^{-s}=\rho^{-s}\omega^{-s}$ the integral is
$\omega^{1-s}(2i)^{-1}\widetilde\Phi(s)$, and $\omega^{1-s}=e^{i\pi(s-1)/4}$.  Absolute
convergence is \Cref{def:rsz_phi_mirror}.  (The Lean statement carries no hypothesis on $s$ at
all: the two integrands are proportional by a constant, so off $\sigma<0$ both sides degenerate to
the junk value $0$ together.)
\end{proof}

\begin{theorem}[The Mellin form of Siegel's identity]
\label{thm:rsz_mellin}
\lean{Zeta.RiemannSiegel.riemannZeta_add_contourIntegral_half_eq}
\leanok
\uses{thm:rsz_mordell, lem:rsz_ray_zeta, lem:rsz_ray_contour, lem:rsz_ray_hankel}
For $\operatorname{Re}s>1$,
\[
  \zeta(s)+I\bigl(s,\tfrac12\bigr)
   \;=\;\frac{(2\pi)^{s}e^{-i\pi(2-s)/4}}{2\,\Gamma(s)}\,\widetilde\Phi(1-s)
   \;=\;\frac{(2\pi)^{s}e^{-i\pi(2-s)/4}}{2\,\Gamma(s)}
        \int_{0}^{\infty}\frac{\rho^{s-1}e^{-\pi\rho^{2}}}{\sin(\pi\rho e^{-i\pi/4})}\,d\rho .
\]
\end{theorem}

\begin{proof}
\leanok
\uses{thm:rsz_mordell, lem:rsz_ray_zeta, lem:rsz_ray_contour, lem:rsz_ray_hankel}
Let first $\sigma<0$.  By \Cref{thm:rsz_mordell},
\[
  N(u)=\frac1{1-e^{2\pi iu}}+\frac{e^{-i\pi u^{2}}}{e^{i\pi u}-e^{-i\pi u}}
  \qquad(u\notin\mathbb{Z}),
\]
since $\frac1{1-e^{2\pi iu}}=\frac{-e^{-i\pi u}}{e^{i\pi u}-e^{-i\pi u}}$.  The ray meets
$\mathbb{Z}$ only at the endpoint $\rho=0$, a null set, so the identity may be multiplied by
$u^{-s}$ and integrated over the ray.  Each of the three resulting integrals converges
absolutely and separately for $\sigma<0$: for the two singular ones because near $u=0$ both
$\frac1{1-e^{2\pi iu}}$ and $\frac{e^{-i\pi u^{2}}}{e^{i\pi u}-e^{-i\pi u}}$ have a simple
pole, so their integrands are $O(\rho^{-\sigma-1})$ and $-\sigma-1>-1$; this is exactly the
hypothesis $\sigma<0$, and it is the only reason for it.  At the other end there is no issue
at all: as $\rho\to\infty$ the geometric term decays like $e^{-\sqrt2\pi\rho}$, the Gaussian
term like $e^{-\pi\rho^{2}}$ and $N$, being their sum, likewise, against the merely polynomial
$\rho^{-\sigma}$.  Applying
\Cref{lem:rsz_ray_contour} on the left and \Cref{lem:rsz_ray_zeta},
\Cref{lem:rsz_ray_hankel} on the right, and abbreviating
$B(s):=(2\pi)^{s-1}e^{i\pi(s-1)/2}\Gamma(1-s)$,
\[
  B(s)\,I\bigl(1-s,\tfrac12\bigr)
   =-B(s)\,\zeta(1-s)+\frac{e^{i\pi(s-1)/4}}{2i}\widetilde\Phi(s),
\]
that is
\begin{equation}
  B(s)\Bigl[\zeta(1-s)+I\bigl(1-s,\tfrac12\bigr)\Bigr]
   =\frac{e^{i\pi(s-1)/4}}{2i}\,\widetilde\Phi(s)\qquad(\sigma<0).
  \label{eq:rsz_mellin_raw}
\end{equation}
Now replace $s$ by $1-s$, so that the hypothesis $\sigma<0$ becomes
$\operatorname{Re}s>1$.  Then $B(1-s)=(2\pi)^{-s}e^{-i\pi s/2}\Gamma(s)$ and the right-hand
side is $e^{-i\pi s/4}\widetilde\Phi(1-s)/2i$, so
\[
  \zeta(s)+I\bigl(s,\tfrac12\bigr)
   =\frac{(2\pi)^{s}e^{i\pi s/2}e^{-i\pi s/4}}{2i\,\Gamma(s)}\widetilde\Phi(1-s)
   =\frac{(2\pi)^{s}e^{i\pi s/4}e^{-i\pi/2}}{2\,\Gamma(s)}\widetilde\Phi(1-s),
\]
using $1/i=e^{-i\pi/2}$, and $e^{i\pi s/4}e^{-i\pi/2}=e^{i\pi(s-2)/4}=e^{-i\pi(2-s)/4}$.
Finally $\widetilde\Phi(1-s)=\int_{0}^{\infty}\rho^{s-1}e^{-\pi\rho^{2}}
/\sin(\pi\rho e^{-i\pi/4})d\rho$ by \Cref{def:rsz_phi_mirror}.
\end{proof}

\subsection{Proof of Siegel's identity}
\label{sec:rsz_proof}

The two objects of \Cref{rem:rsz_general} are now given names for the formalization.

\begin{definition}[The mirror contour integral and the functional-equation factor]
\label{def:rsz_mirror}
\lean{Zeta.RiemannSiegel.mirrorContourIntegral, Complex.chi}
\leanok
\uses{def:rs_integral}
For $s\in\mathbb{C}$ and $a\in\tfrac12+\mathbb{Z}_{\ge0}$ set
\[
  \widetilde I(s,a):=\overline{I\bigl(\overline{s},a\bigr)},
  \qquad
  \chi(s):=\pi^{s-\frac12}\,
     \frac{\Gamma\bigl(\frac{1-s}2\bigr)}{\Gamma\bigl(\frac s2\bigr)} .
\]
The first is a literal transcription of \Cref{rem:rsz_general}; in particular $\widetilde I$
needs no second contour and no second branch of the power, being defined by conjugation from
$I$.

The second differs from the formula $2^{s}\pi^{s-1}\sin(\pi s/2)\Gamma(1-s)$ of
\Cref{rem:rsz_general}, and the quotient form is the one formalized.  The two agree as
meromorphic functions
(\Cref{lem:rsz_chi}), but mathlib's \texttt{Complex.Gamma} takes the junk value $0$ at its
poles, so the product form evaluates to $\bigl(\cdots\bigr)\cdot0\cdot0=0$ at every even
integer $s\ge2$, where the true value is not $0$ (for instance $\chi(2)=-2\pi^{2}$).  Under
the product form \Cref{thm:rsz_general} would assert $\zeta(2)+I(2,\tfrac12)=0$, which is
false: $\widetilde I(-1,\tfrac12)=\widetilde\Phi(-1)=\tfrac18+\tfrac i{4\pi}$ by
\Cref{lem:rsz_mirror}, and $\zeta(2)+I(2,\tfrac12)=-\chi(2)\widetilde I(-1,\tfrac12)
=\tfrac{\pi^{2}}4+\tfrac{\pi i}2$.  The quotient form has no such defect (in Lean it is taken in Deligne's normalisation
$\chi(s)=\Gamma_{\mathbb{R}}(1-s)/\Gamma_{\mathbb{R}}(s)$, which is the same function with no
exclusions, the displayed form being \texttt{Complex.chi\_eq\_cpow\_mul\_Gamma\_div\_Gamma}):
$1/\Gamma$ is entire, so the only junk values are at the
poles of $\Gamma\bigl(\frac{1-s}2\bigr)$, namely $s=1,3,5,\dots$, which is exactly the set
excluded in \Cref{thm:rsz_general}; and at $s=0,-2,-4,\dots$ the junk value $1/\Gamma(0)=0$
of the denominator agrees with the true value $\chi(s)=0$.
\end{definition}

\begin{lemma}[The mirror integral in Hankel form]
\label{lem:rsz_mirror}
\lean{Zeta.RiemannSiegel.mirrorContourIntegral_one_sub_eq}
\leanok
\uses{lem:rsz_hankel, def:rsz_phi_mirror, def:rsz_mirror}
For $\operatorname{Re}s>1$,
\[
  \widetilde I\bigl(1-s,\tfrac12\bigr)
   \;=\;-\,e^{-i\pi(2-s)/4}\cos\Bigl(\frac{\pi s}{2}\Bigr)\,\widetilde\Phi(1-s).
\]
\end{lemma}

\begin{proof}
\leanok
\uses{lem:rsz_hankel, def:rsz_phi_mirror}
By definition $\widetilde I(1-s,\tfrac12)=\overline{I(1-\overline{s},\tfrac12)}$, and
$\operatorname{Re}(1-\overline s)=1-\operatorname{Re}s<0$, so \Cref{lem:rsz_hankel} applies at
$z=1-\overline s$ and gives
$I(z,\tfrac12)=-e^{i\pi(1+z)/4}\sin(\pi z/2)\Phi(z)$.  Conjugating,
$\overline{e^{i\pi(2-\overline s)/4}}=e^{-i\pi(2-s)/4}$,
$\overline{\sin(\pi(1-\overline s)/2)}=\sin(\pi(1-s)/2)=\cos(\pi s/2)$, and
$\overline{\Phi(1-\overline s)}=\widetilde\Phi(1-s)$ by \Cref{def:rsz_phi_mirror}.
\end{proof}

\begin{lemma}[Two identities for $\chi$]
\label{lem:rsz_chi}
\lean{Complex.chi_mul_cos, Complex.chi_eq_two_cpow_mul_cpow_mul_sin_mul_Gamma}
\leanok
\uses{def:rsz_mirror}
Let $\chi$ be as in \Cref{def:rsz_mirror}.  Then
\begin{equation}
  \chi(s)\cos\Bigl(\frac{\pi s}2\Bigr)=\frac{(2\pi)^{s}}{2\,\Gamma(s)}
    \qquad\bigl(s\notin\{1,3,5,\dots\}\bigr),
  \label{eq:rsz_chi_cos}
\end{equation}
\begin{equation}
  \chi(s)=2^{s}\pi^{s-1}\sin\Bigl(\frac{\pi s}2\Bigr)\Gamma(1-s)
    \qquad\bigl(s\notin\{2,4,6,\dots\}\bigr).
  \label{eq:rsz_chi_prod}
\end{equation}
Each exclusion is necessary under mathlib's junk-value convention: at $s=3,5,\dots$ the
left side of \eqref{eq:rsz_chi_cos} is $0$ (both factors
are) while the right side is not, and at $s=2,4,6,\dots$ the right side of
\eqref{eq:rsz_chi_prod} is $0$ while the left side is not (for instance both sides of
\eqref{eq:rsz_chi_prod} at $s=2$ are $-2\pi^{2}$ and $0$).  The two excluded sets are
different, and both are sharp: at the \emph{odd} positive integers, where $\Gamma(1-s)$ is
also singular, both sides of \eqref{eq:rsz_chi_prod} do vanish, so those points need not be
excluded and are not.
\end{lemma}

\begin{proof}
\leanok
For \eqref{eq:rsz_chi_cos}, first dispose of $s\in\{0,-1,-2,\dots\}$: there $\Gamma(s)$ is
singular, so the right side is $(2\pi)^{s}/(2\cdot0)=0$, and the left side vanishes too:
at $s=0,-2,-4,\dots$ because $\Gamma(s/2)$ is singular and $\chi$ has it in the
\emph{denominator}, at $s=-1,-3,-5,\dots$ because $\cos(\pi s/2)=0$.  Off that set both
$\Gamma(s)$ and $\Gamma(s/2)$ are nonzero, so one may clear denominators: multiplying by
$2\,\Gamma(s)\,\Gamma(s/2)\,\pi^{\frac12-s}$ and using $(2\pi)^{s}\pi^{\frac12-s}
=2^{s}\sqrt\pi$, the claim becomes
\[
  2\,\Gamma(s)\,\Gamma\Bigl(\tfrac{1-s}2\Bigr)\cos\Bigl(\tfrac{\pi s}2\Bigr)
  \;=\;2^{s}\sqrt\pi\;\Gamma\Bigl(\tfrac s2\Bigr).
\]
Legendre duplication (mathlib's
\texttt{Complex.Gamma\_mul\_Gamma\_add\_half}, $\Gamma(z)\Gamma(z+\tfrac12)
=2^{1-2z}\sqrt\pi\,\Gamma(2z)$) at $z=\tfrac s2$ gives
$\Gamma(s)=2^{s-1}\Gamma\bigl(\tfrac s2\bigr)\Gamma\bigl(\tfrac{s+1}2\bigr)/\sqrt\pi$, so the
left side is
$2^{s}\pi^{-1/2}\Gamma\bigl(\tfrac s2\bigr)\cdot
\Gamma\bigl(\tfrac{s+1}2\bigr)\Gamma\bigl(\tfrac{1-s}2\bigr)\cos\bigl(\tfrac{\pi s}2\bigr)$,
and reflection (\texttt{Complex.Gamma\_mul\_Gamma\_one\_sub}) at $z=\tfrac{1+s}2$, whose
partner is $1-z=\tfrac{1-s}2$, gives
$\Gamma\bigl(\tfrac{s+1}2\bigr)\Gamma\bigl(\tfrac{1-s}2\bigr)
=\pi/\sin\bigl(\tfrac{\pi(1+s)}2\bigr)=\pi/\cos\bigl(\tfrac{\pi s}2\bigr)$.  The two cosines
cancel and $2^{s}\pi^{-1/2}\Gamma(\tfrac s2)\cdot\pi=2^{s}\sqrt\pi\,\Gamma(\tfrac s2)$, as
required.  The only division here is by $\cos(\pi s/2)$, in the reflection formula as
rearranged; it is legitimate exactly off $s\in\{1,3,5,\dots\}$, and that is where the
excluded set of \eqref{eq:rsz_chi_cos} comes from.

For \eqref{eq:rsz_chi_prod}, duplication at $z=\tfrac{1-s}2$ gives
$\Gamma\bigl(\tfrac{1-s}2\bigr)\Gamma\bigl(1-\tfrac s2\bigr)=2^{s}\sqrt\pi\,\Gamma(1-s)$,
and reflection at $z=\tfrac s2$ gives $\sin(\pi s/2)\Gamma\bigl(1-\tfrac s2\bigr)
=\pi/\Gamma\bigl(\tfrac s2\bigr)$, a division by $\Gamma(s/2)$, legitimate because
$\Gamma$ vanishes nowhere and, at the poles $s=0,-2,-4,\dots$ of $\Gamma(s/2)$, both sides
of \eqref{eq:rsz_chi_prod} are $0$ by the same inspection as above.  Whence
\[
  2^{s}\pi^{s-1}\sin\Bigl(\tfrac{\pi s}2\Bigr)\Gamma(1-s)
   =\pi^{s-\frac32}\sin\Bigl(\tfrac{\pi s}2\Bigr)
     \Gamma\Bigl(\tfrac{1-s}2\Bigr)\Gamma\Bigl(1-\tfrac s2\Bigr)
   =\pi^{s-\frac12}\frac{\Gamma\bigl(\tfrac{1-s}2\bigr)}{\Gamma\bigl(\tfrac s2\bigr)}
   =\chi(s).
\]
The division by $\sin(\pi s/2)$ that this rearrangement conceals is legitimate exactly off the
even integers, and the nonpositive ones are the case $\Gamma(s/2)$ singular already disposed
of; that is where the excluded set of \eqref{eq:rsz_chi_prod} comes from.

Both identities are formalized through Deligne's archimedean factor
$\Gamma_{\mathbb{R}}(s)=\pi^{-s/2}\Gamma(s/2)$ of mathlib's \texttt{Gamma/Deligne.lean}, which
packages duplication and reflection together.  One has
$\chi(s)=\Gamma_{\mathbb{R}}(1-s)/\Gamma_{\mathbb{R}}(s)$ identically (with no exclusions,
since $\pi^{w}$ never vanishes), and mathlib's reflection formula
$\Gamma_{\mathbb{R}}(s)/\Gamma_{\mathbb{R}}(1-s)=\Gamma_{\mathbb{C}}(s)\cos(\pi s/2)$, valid
for $s\notin\{-1,-3,-5,\dots\}$, is \eqref{eq:rsz_chi_cos} on inversion and, read at $1-s$
instead of at $s$, is \eqref{eq:rsz_chi_prod}.  Reading it at $1-s$ is exactly what produces
the excluded set $\{2,4,6,\dots\}$ rather than $\{1,2,3,\dots\}$.
\end{proof}

\begin{remark}[The functional equation is not used]
\label{rem:rsz_chi_fe}
\uses{lem:rsz_chi, rem:rsz_general}
\Cref{rem:rsz_general} characterises $\chi$ by $\zeta(s)=\chi(s)\zeta(1-s)$, and
\eqref{eq:rsz_chi_prod} identifies \Cref{def:rsz_mirror} with the classical formula, so the
$\chi$ of this section is the classical one.  But the functional equation itself is nowhere
used below: \Cref{prop:rsz_general_half} uses only \eqref{eq:rsz_chi_cos} and
\Cref{lem:rsz_chi_theta} uses only \Cref{def:rsz_mirror}.  This is worth recording because
mathlib's \texttt{riemannZeta\_one\_sub} carries the hypotheses $s\ne-n$ $(n\in\mathbb{N})$
and $s\ne1$, which after the substitution $s\mapsto1-s$ exclude $s=0$ and
$s\in\{1,2,3,\dots\}$; none of that has to be managed here.  It is also what makes
\Cref{thm:rsz_general} a second proof of the functional equation, as
\cite{Siegel1932} observes: on $\sigma=\tfrac12$ the two contour terms are complex
conjugates, so $\pi^{-s/2}\Gamma(s/2)\zeta(s)$ is real there.
\end{remark}

\begin{proposition}[Siegel's identity on $\operatorname{Re}s>1$]
\label{prop:rsz_general_half}
\lean{Zeta.RiemannSiegel.riemannZeta_add_contourIntegral_add_chi_mul_mirror_eq_zero_of_one_lt}
\leanok
\uses{thm:rsz_mellin, lem:rsz_mirror, lem:rsz_chi, def:rsz_mirror}
For $\operatorname{Re}s>1$ with $s\notin\{3,5,7,\dots\}$,
\[
  \zeta(s)+I\bigl(s,\tfrac12\bigr)+\chi(s)\,\widetilde I\bigl(1-s,\tfrac12\bigr)=0 .
\]
\end{proposition}

\begin{proof}
\uses{thm:rsz_mellin, lem:rsz_mirror, lem:rsz_chi}
\leanok
By \Cref{lem:rsz_mirror} and then \eqref{eq:rsz_chi_cos},
\[
  -\chi(s)\widetilde I\bigl(1-s,\tfrac12\bigr)
   =e^{-i\pi(2-s)/4}\,\chi(s)\cos\Bigl(\frac{\pi s}2\Bigr)\widetilde\Phi(1-s)
   =\frac{(2\pi)^{s}e^{-i\pi(2-s)/4}}{2\Gamma(s)}\widetilde\Phi(1-s),
\]
which is $\zeta(s)+I(s,\tfrac12)$ by \Cref{thm:rsz_mellin}.  The second step is the only one
that needs $s\notin\{3,5,7,\dots\}$: at those points $\chi$ has a simple pole, while
$\cos(\pi s/2)$ and hence $\widetilde I(1-s,\tfrac12)$ vanish, so the product
$\chi(s)\widetilde I(1-s,\tfrac12)$ is of the form $\infty\cdot0$ and the statement as
displayed has no content there.  These points stay excluded in \Cref{thm:rsz_general}; see
\Cref{rem:rsz_general}, which explains why the natural domain of
\eqref{eq:rsz_siegel_general} is $\mathbb{C}\setminus\{1,3,5,\dots\}$ and gives a pole-free
rewriting.  $\Gamma(s)$ is finite and nonvanishing throughout $\operatorname{Re}s>1$.
\end{proof}

\begin{lemma}[$I(\cdot,a)$ is entire]
\label{lem:rsz_entire}
\lean{Zeta.RiemannSiegel.differentiable_contourIntegral}
\leanok
\uses{def:rs_integral, prop:rs_integrable, lem:rs_power_modulus, lem:rs_rotate}
For each $a\in\tfrac12+\mathbb{Z}_{\ge0}$ the map $s\mapsto I(s,a)$ is entire, and hence so
is $s\mapsto\widetilde I(1-s,a)=\overline{I(1-\overline s,a)}$.

In Lean the hypothesis is taken in the same form as for
\texttt{integrable\_contourIntegrand}, namely a lower bound $c>0$ for the kernel denominator
along the contour, with the two corollaries
\texttt{differentiable\_contourIntegral\_of\_eq\_add\_half} (half-integer $a$, the statement
above) and \texttt{differentiable\_contourIntegral\_of\_le\_half} ($0<a\le\tfrac12$, which is
the case $a=\tfrac12$ that \Cref{thm:rsz_general} uses).  The mirror half is
\texttt{differentiable\_mirrorContourIntegral\_one\_sub}.
\end{lemma}

\begin{proof}
\uses{prop:rs_integrable, lem:rs_power_modulus, lem:rs_rotate}
\leanok
Differentiation under the integral sign.  For each $w$ the integrand
$s\mapsto L(a)(w)^{-s}G(L(a)(w))$ is entire, with $s$-derivative
$-\operatorname{Log}L(a)(w)\cdot L(a)(w)^{-s}G(L(a)(w))$.  On a compact set of $s$-values,
$\lvert\sigma\rvert\le A$ and $\lvert\tau\rvert\le A$, the bounds of
\Cref{prop:rs_integrable} give
$\lvert L(a)(w)^{-s}\rvert\le\lvert L(a)(w)\rvert^{-\sigma}e^{3\pi A/4}$ with
$a/\sqrt2\le\lvert L(a)(w)\rvert\le a+2\lvert w\rvert$ (\Cref{lem:rs_rotate}), so the
integrand and its $s$-derivative are both dominated by
$C_{A}\bigl(1+\lvert w\rvert\bigr)^{A}\bigl(1+\bigl\lvert\log(a+2\lvert
w\rvert)\bigr\rvert+\lvert\log(a/\sqrt2)\rvert\bigr)e^{-\pi w^{2}}$, which is integrable.
Mathlib's \texttt{hasDerivAt\_integral\_of\_dominated\_loc\_of\_deriv\_le} then gives
differentiability at every $s$.  Conjugation reverses no analyticity here because it is
applied twice: $s\mapsto\overline{I(1-\overline s,a)}$ is a composition of two conjugations
with an entire map.
\end{proof}

\begin{theorem}[Siegel's identity, general form]
\label{thm:rsz_general}
\lean{Zeta.RiemannSiegel.riemannZeta_add_contourIntegral_add_chi_mul_mirror_eq_zero}
\leanok
\uses{prop:rsz_general_half, lem:rsz_entire, def:rsz_mirror}
Identity \eqref{eq:rsz_siegel_general} holds for every
$s\in\mathbb{C}\setminus\{1,3,5,7,\dots\}$:
\[
  \zeta(s)+I\bigl(s,\tfrac12\bigr)+\chi(s)\,\widetilde I\bigl(1-s,\tfrac12\bigr)=0 .
\]
\end{theorem}

\begin{proof}
\uses{prop:rsz_general_half, lem:rsz_entire}
\leanok
Let $U:=\mathbb{C}\setminus\{1,3,5,7,\dots\}$.  This is the complement of a discrete set in
$\mathbb{C}$, hence open and connected.  On $U$ the function
$F(s):=\zeta(s)+I(s,\tfrac12)+\chi(s)\widetilde I(1-s,\tfrac12)$ is holomorphic: $\zeta$ is
holomorphic off $s=1$; $\chi$ is holomorphic on $U$ because, in the quotient form of
\Cref{def:rsz_mirror}, $1/\Gamma$ is \emph{entire} and the only poles of
$\Gamma\bigl(\frac{1-s}2\bigr)$ are at $\frac{1-s}2\in\{0,-1,-2,\dots\}$, i.e.\ at
$s\in\{1,3,5,\dots\}$ (the product form would need an argument that the poles of
$\Gamma(1-s)$ at the even integers are cancelled by the zeros of $\sin(\pi s/2)$); and both
contour integrals are entire
by \Cref{lem:rsz_entire}.  By \Cref{prop:rsz_general_half}, $F$ vanishes on
$\{\operatorname{Re}s>1\}\cap U$, which is open and nonempty.  By the identity theorem
$F\equiv0$ on $U$.
\end{proof}

\begin{lemma}[$\chi$ on the critical line]
\label{lem:rsz_chi_theta}
\lean{Complex.chi_half_add_mul_I}
\leanok
\uses{def:rsz_mirror, def:nf_theta}
For every real $t\ge0$, $\chi\bigl(\tfrac12+it\bigr)=e^{-2i\theta(t)}$.
\end{lemma}

\begin{proof}
\uses{def:rsz_mirror, def:nf_theta}
\leanok
Directly from \Cref{def:rsz_mirror} at $s=\tfrac12+it$; the quotient form makes the lemma a
computation with no $\Gamma$ identity at all:
\[
  \chi\bigl(\tfrac12+it\bigr)=\pi^{it}\,
    \frac{\Gamma\bigl(\tfrac14-\tfrac{it}2\bigr)}{\Gamma\bigl(\tfrac14+\tfrac{it}2\bigr)} ,
\]
because $\tfrac12\bigl(1-(\tfrac12+it)\bigr)=\tfrac14-\tfrac{it}2$ and
$\tfrac12\bigl(\tfrac12+it\bigr)=\tfrac14+\tfrac{it}2$.  The restriction $t\ge0$ is inherited
from \Cref{def:nf_theta}, which defines $\theta$ by a continuous branch along the segment
$[0,T]$; nothing in the argument needs it.
Write $z:=\log\Gamma\bigl(\tfrac14+\tfrac{it}2\bigr)$ along the branch of
\Cref{def:nf_theta}.  Since $\Gamma$ is real on $(0,\infty)$,
$\Gamma\bigl(\tfrac14-\tfrac{it}2\bigr)=\overline{\Gamma\bigl(\tfrac14+\tfrac{it}2\bigr)}
=e^{\overline z}$, so the quotient is $e^{\overline z-z}=e^{-2i\operatorname{Im}z}$.  By
\Cref{def:nf_theta}, $\operatorname{Im}z=\theta(t)+\tfrac t2\log\pi$, whence the quotient is
$e^{-2i\theta(t)}\pi^{-it}$, and the factor $\pi^{it}$ cancels.
\end{proof}

\begin{theorem}[Siegel's identity, critical line]
\label{thm:rsz_siegel}
\lean{Zeta.RiemannSiegel.riemannZeta_eq_neg_contourIntegral_sub}
\leanok
\uses{def:rs_integral, def:nf_theta, def:rsz_mirror, thm:rsz_general, lem:rsz_chi_theta}
For every $t>0$, with $s=\tfrac12+it$,
\[
  \zeta(s)\;=\;-\,I\bigl(s,\tfrac12\bigr)\;-\;e^{-2i\theta(t)}\,\overline{I\bigl(s,\tfrac12\bigr)} .
\]
\end{theorem}

\begin{proof}
\uses{thm:rsz_general, lem:rsz_chi_theta}
\leanok
The point $s=\tfrac12+it$ lies in the set $U$ of \Cref{thm:rsz_general}, since its real part
is $\tfrac12$.  There $1-s=\tfrac12-it=\overline s$, so
$\overline{1-s}=s$ and therefore
\[
  \widetilde I\bigl(1-s,\tfrac12\bigr)=\overline{I\bigl(\overline{1-s},\tfrac12\bigr)}
    =\overline{I\bigl(s,\tfrac12\bigr)} .
\]
Substituting this and \Cref{lem:rsz_chi_theta} into \Cref{thm:rsz_general} and solving for
$\zeta(s)$ gives the claim.
\end{proof}

\begin{remark}[Two routes that do not work, and why]
\label{rem:rsz_routes}
\uses{lem:rsz_hankel, lem:rsz_expansion, lem:rsz_gauss_integral, rem:rsz_hankel_mirror,
thm:rsz_mellin}
Both of the obvious attacks on the crossing fail.
\begin{enumerate}
  \item \emph{Rotating the Hankel ray.}  \Cref{lem:rsz_hankel} represents $I(s,\tfrac12)$ by
        an integral along the ray $e^{i\pi/4}(0,\infty)$ and
        \Cref{rem:rsz_hankel_mirror} represents the mirror half by one along
        $e^{-i\pi/4}(0,\infty)$, so it is tempting to rotate one into the other, collecting
        the residues at the positive integers in between.  The residue bookkeeping is
        correct (the residue at $n$ is $n^{s-1}/\pi$, because $e^{-i\pi n^{2}}=(-1)^{n}$
        cancels the $(-1)^{n}/\pi$ of $1/\sin$, and the sum is $2i\zeta(1-s)$), but the
        rotation is meaningless: on $z=te^{i\pi/4}$ one has $z^{2}=it^{2}$, so the mirror
        Gaussian $e^{-i\pi z^{2}}$ becomes $e^{+\pi t^{2}}$ and the integral diverges.  The two
        Hankel representations
        live on rays that can never be deformed into one another, whatever the truncation.
  \item \emph{Term-by-term expansion of the kernel.}  Inserting \Cref{lem:rsz_expansion} into
        $I(s,a)$ and integrating term by term against $x^{-s}$ does not close either.  The
        expansion is valid only on a half-plane, so the contour must be split at $w=0$ and
        each term becomes an \emph{incomplete} Gaussian moment, not the closed value
        $e^{i\pi/4}$ of \Cref{lem:rsz_gauss_integral}; and the individual terms are enormous
        (by \Cref{rem:rsz_cancellation} the $k=0$ term alone has modulus
        $e^{\pi a(a-1)/2}$), so no rearrangement can be justified by moduli.
\end{enumerate}
\Cref{thm:rsz_mellin} introduces a \emph{second} variable.  The
Mellin transform in $u$ never moves the contour in $x$, and the only convergence that has to
be checked is the one supplied by \Cref{lem:rs_rotate}: $\operatorname{Re}(e^{i\pi/4}x)$ is
constant on $L(a)$.  The entire-function identity
$G(x)+\widetilde G(x)=\sin(\pi x^{2})/\sin(\pi x)$ of \Cref{rem:rs_mirror} plays no role: it
cannot, since neither $G$ nor $\widetilde G$ is integrable against $x^{-s}$ on the other's
contour.
\end{remark}

\section{The order-zero formula}
\label{sec:rsz_formula}

\begin{theorem}[Order-zero Riemann--Siegel formula, exact form]
\label{thm:rsz_formula}
\lean{Zeta.RiemannSiegel.hardyZ_eq_mainSum_sub}
\leanok
\uses{thm:rsz_reduction, thm:rsz_siegel, def:hz_Z, def:rsz_main_sum}
For every $t>0$ and every $\nu\in\mathbb{N}$, with $s=\tfrac12+it$,
\[
  Z(t)\;=\;\mathcal{S}(t,\nu)\;-\;2\operatorname{Re}
     \Bigl(e^{i\theta(t)}\,I\bigl(s,\nu+\tfrac12\bigr)\Bigr).
\]
\end{theorem}

\begin{proof}
\uses{thm:rsz_siegel, thm:rsz_reduction}
\leanok
Abbreviate $\theta=\theta(t)$ and $J=I(s,\tfrac12)$.  Multiplying \Cref{thm:rsz_siegel} by
$e^{i\theta}$ gives
\[
  e^{i\theta}\zeta(s)=-e^{i\theta}J-e^{-i\theta}\overline{J}
   =-\bigl(e^{i\theta}J+\overline{e^{i\theta}J}\bigr)
   =-2\operatorname{Re}\bigl(e^{i\theta}J\bigr),
\]
a real number.  By \Cref{def:hz_Z}, $Z(t)=\operatorname{Re}\bigl(e^{i\theta}\zeta(s)\bigr)$,
so $Z(t)=-2\operatorname{Re}(e^{i\theta}J)$.  Adding \Cref{thm:rsz_reduction} gives the
claim.
\end{proof}

\begin{remark}[Axioms]
\label{rem:rsz_conditional}
\uses{thm:rsz_formula, cor:rsz_bound, thm:rsz_siegel}
\Cref{thm:rsz_formula} and \Cref{cor:rsz_bound} take Siegel's identity from
\Cref{thm:rsz_siegel}, under the hypothesis $t>0$ that it requires.  \Cref{cha:rs_zero} and
\Cref{cha:rs_first} depend on no axioms beyond
\texttt{propext}, \texttt{Classical.choice} and \texttt{Quot.sound}.
\end{remark}

\begin{remark}[Realness of $Z$]
\label{rem:rsz_realness}
\uses{thm:rsz_siegel}
The computation in the proof shows that $e^{i\theta(t)}\zeta(\tfrac12+it)$ is real, i.e.\
that $Z(t)=e^{i\theta(t)}\zeta(\tfrac12+it)$ and not merely its real part.  The same fact
appears inside the proof of
\texttt{hardyZ\_eq\_zero\_iff\_riemannZeta}, via the realness of $\xi$ on the critical line.
\Cref{thm:rsz_formula} uses only the definition of $Z$ as a real part.
\end{remark}

\begin{corollary}[Order-zero bound]
\label{cor:rsz_bound}
\lean{Zeta.RiemannSiegel.abs_hardyZ_sub_mainSum_le_hundredNinety, Zeta.RiemannSiegel.abs_hardyZ_sub_mainSum_le_threeHundredTwo}
\leanok
\uses{thm:rsz_formula, cor:rs_remainder_explicit}
Let $t\ge32\pi$, $T=t/2\pi$ and $\nu=\lfloor\sqrt T\rfloor$.  Then
\[
  \bigl\lvert Z(t)-\mathcal{S}(t,\nu)\bigr\rvert\;\le\;190\,T^{-1/4}\;\le\;302\,t^{-1/4}.
\]
\end{corollary}

\begin{proof}
\uses{thm:rsz_formula, cor:rs_remainder_explicit}
\leanok
By \Cref{thm:rsz_formula} the difference equals
$-2\operatorname{Re}\bigl(e^{i\theta(t)}I(s,\nu+\tfrac12)\bigr)$, whose modulus is at most
$2\lvert I(s,\nu+\tfrac12)\rvert$ because $\lvert e^{i\theta(t)}\rvert=1$.  Since
$\nu\le\sqrt T<\nu+1$, \Cref{cor:rs_remainder_explicit} applies with $a=\nu+\tfrac12$ and
gives $\lvert I(s,\nu+\tfrac12)\rvert\le95\,T^{-1/4}\le151\,t^{-1/4}$.  Double both
constants.
\end{proof}

\begin{remark}[Cost and accuracy]
\label{rem:rsz_next}
\uses{cor:rsz_bound}
\Cref{cor:rsz_bound} evaluates $Z(t)$ to within $190\,T^{-1/4}$ using
$\lfloor\sqrt{T}\rfloor$ terms, against the $\Theta(t)$ terms per point of the
Euler--Maclaurin evaluation of \Cref{cha:hardy_z}.  This alone does not certify zeros: the
error $190\,T^{-1/4}$ exceeds the typical size of $Z$ until $T$ is astronomically large.  The
following chapters make two independent improvements.
\begin{itemize}
  \item The constant.  As \Cref{rem:rs_constant} explains, essentially all of the factor
        $95$ is the price $e^{B^{2}/4\pi}\le37.12$ of crossing the real axis at
        $\nu+\tfrac12$ rather than at the saddle $\sqrt T$.  Recovering it produces the
        classical zeroth correction term
        $(-1)^{\nu-1}T^{-1/4}\Psi(p)$, $\Psi(p)=\cos(2\pi(p^{2}-p-\tfrac1{16}))/\cos(2\pi p)$,
        with $p=\sqrt T-\nu$.
  \item The order.  Subtracting $\Psi$ leaves $O(T^{-3/4})$, and the general saddle-point
        expansion leaves $O(T^{-(2K+1)/4})$ after $K$ correction terms; see
        \cite{Gabcke} and \cite{AriasDeReyna2011} for explicit constants.
\end{itemize}
Both improvements keep \Cref{thm:rsz_reduction} and \Cref{thm:rsz_siegel} untouched: they
refine only the treatment of $I(s,\nu+\tfrac12)$, which \Cref{thm:rsz_formula} isolates as
the exact error.
\end{remark}
